\section{Exactly summable commutator relations}\label{sec:special}
The universal series is infinite, but algebraic relations can remove almost
all of it. The proofs below establish \emph{global} exponential identities,
not merely small-neighborhood Taylor expansions. For Banach algebras all
exponentials exist; for finite-dimensional Lie groups the same proofs use
one-parameter subgroups and their right-trivialized logarithmic derivatives
(\cref{prop:intrinsicdexp}), together with uniqueness of solutions of the
ordinary differential equation $g'=V(t)^R(g)$ on the group. Complex
scalar parameters are understood in a complex Banach algebra or complex Lie
group; in a real Lie group all parameters and linear combinations must be
real.

\subsection{Commuting elements and central commutators}
\begin{proposition}[Commuting case]\label{prop:commuting}
If $[X,Y]=0$, then $Z(X,Y)=X+Y$ formally, and $e^Xe^Y=e^{X+Y}$ for all
bounded $X,Y$ without a smallness hypothesis.
\end{proposition}
\begin{proof}
For commuting elements the binomial theorem holds, and the Cauchy product
of the two absolutely convergent series gives
$e^Xe^Y=\sum_{r,s}X^rY^s/(r!s!)=\sum_n\frac1{n!}\sum_{r+s=n}\binom nrX^rY^s
=\sum_n(X+Y)^n/n!=e^{X+Y}$. Formally, the same computation is valid
coefficientwise, and then $Z(X,Y)=\log e^{X+Y}=X+Y$ by \cref{lem:explog};
this is also the content of \cref{thm:formalbch}, since every $Z_n$ with
$n\geq2$ lies in the ideal generated by $[X,Y]$.
\end{proof}

\begin{theorem}[Central-commutator identity]\label{thm:central}
Suppose $C=[X,Y]$ commutes with both $X$ and $Y$. Then, for every scalar
parameter $t$,
\begin{align}
 e^{tX}Ye^{-tX}&=Y+tC,\label{eq:centralconj}\\
 e^{tX}e^{tY}&=\exp\bigl(t(X+Y)+\tfrac12t^2C\bigr).\label{eq:centralt}
\end{align}
In particular,
\begin{equation}\label{eq:centralbch}
 e^Xe^Y=e^{X+Y+C/2},\qquad
 e^{X+Y}=e^Xe^Ye^{-C/2},\qquad
 e^{X/2}e^Ye^{X/2}=e^{X+Y},\qquad
 e^Xe^Ye^{-X}e^{-Y}=e^{C}.
\end{equation}
These are formal identities, global identities for bounded elements of a
Banach algebra, and global identities for elements of a finite-dimensional
Lie group whose Lie algebra elements satisfy the stated bracket relations.
\end{theorem}
\begin{proof}
In the Campbell series \eqref{eq:campbell}, $\ad_X^2Y=[X,C]=0$, so all
terms after the first commutator vanish, giving \eqref{eq:centralconj}.
Let $g(t)=e^{tX}e^{tY}$. Its right logarithmic derivative is
\[
 g'(t)g(t)^{-1}=X+e^{tX}Ye^{-tX}=X+Y+tC,
\]
so $g'=(X+Y+tC)g$ with $g(0)=I$. Let $K(t)=t(X+Y)+\frac12t^2C$. Since $C$
commutes with $X+Y$, $K(t)$ commutes with $K'(t)=X+Y+tC$, and therefore
$\ad_{K}K'=0$ and \eqref{eq:rightdexp} gives
$\frac{\dd}{\dd t}e^{K(t)}=\phi(-\ad_K)K'\,e^{K}=K'e^{K}=(X+Y+tC)e^{K(t)}$,
with $e^{K(0)}=I$. Both $g$ and $e^{K}$ solve $h'=(X+Y+tC)h$, $h(0)=I$;
the difference $\Delta$ of two solutions satisfies
$\Delta'=(X+Y+tC)\Delta$, $\Delta(0)=0$, hence $\Delta\equiv0$ by
uniqueness for linear equations with continuous bounded coefficients
(Gr\"onwall), or by comparing coefficients of $t$ formally. This proves
\eqref{eq:centralt}. In a Lie group the same argument uses the
right-trivialized derivative and uniqueness of integral curves of the
time-dependent right-invariant field $t\mapsto(X+Y+tC)^R$.

Setting $t=1$ gives the first identity of \eqref{eq:centralbch}; since
$C/2$ commutes with $X+Y+C/2$, \cref{prop:commuting} gives
$e^{X+Y+C/2}e^{-C/2}=e^{X+Y}$, the second identity. For the symmetric
identity apply the first identity to the pair $(X/2,Y)$, whose commutator
is $C/2$, and then to the pair $(X/2+Y+C/4,\,X/2)$, whose commutator is
$\frac12[Y,X]=-C/2$ (all other brackets vanish and $C$ is central):
\[
 e^{X/2}e^Ye^{X/2}=e^{X/2+Y+C/4}e^{X/2}=e^{X/2+Y+C/4+X/2-C/4}=e^{X+Y}.
\]
Finally, by \eqref{eq:centralconj} and multiplicativity of conjugation,
$e^Xe^Ye^{-X}=e^{Y+C}$, and since $Y$ and $C$ commute,
$e^{Y+C}e^{-Y}=e^{C}$.
\end{proof}
Centrality is needed only relative to the Lie subalgebra generated by $X$
and $Y$, not relative to a larger ambient algebra. This proof is global and
does not infer convergence from the formal truncation alone; for unbounded
operators a common invariant domain and integrability hypotheses are still
necessary, as \cref{sec:quantum} shows.

\begin{corollary}[Several factors with central pairwise commutators]\label{cor:centralmany}
If all commutators $[X_i,X_j]$ commute with every $X_k$, then
\begin{equation}\label{eq:centralmany}
 \prod_{j=1}^{m}e^{X_j}
 =\exp\Bigl(\sum_{j=1}^{m}X_j+\frac12\sum_{1\leq i<j\leq m}[X_i,X_j]\Bigr),
\end{equation}
the product being ordered by increasing $j$.
\end{corollary}
\begin{proof}
Induct on $m$; the case $m=1$ is trivial. Let $E_{m-1}=\sum_{j<m}X_j
+\frac12\sum_{i<j<m}[X_i,X_j]$, so that $\prod_{j<m}e^{X_j}=e^{E_{m-1}}$ by
the induction hypothesis. The commutator $[E_{m-1},X_m]=\sum_{j<m}[X_j,X_m]$,
because the central part of $E_{m-1}$ commutes with $X_m$; it is a sum of
central elements, hence commutes with $E_{m-1}$ and $X_m$. \Cref{thm:central}
applied to $(E_{m-1},X_m)$ gives
$e^{E_{m-1}}e^{X_m}=\exp\bigl(E_{m-1}+X_m+\frac12\sum_{j<m}[X_j,X_m]\bigr)$,
which is \eqref{eq:centralmany} for $m$ factors.
\end{proof}
This is the all-factors form behind the Weyl and displacement
multiplication laws of \cref{sec:quantum}. In the Heisenberg algebra all
higher BCH terms vanish because every Lie monomial of degree at least three
contains a bracket with the central element $C$, in agreement with
\cref{thm:nilpotent}.

\subsection{An eigenvector of the adjoint action}
\begin{theorem}[The relation $[X,Y]=sY$]\label{thm:eigen}
Suppose $[X,Y]=sY$ with $s$ a scalar. Put
\[
 q_s(t)=\int_0^te^{-su}\,\dd u=\begin{cases}(1-e^{-st})/s,&s\neq0,\\ t,&s=0,\end{cases}
 \qquad q_s(1)=\phi(s).
\]
Then for every scalar $c$ and every $t$,
\begin{equation}\label{eq:eigenfactor}
 e^{t(X+cY)}=e^{tX}e^{c\,q_s(t)Y},\qquad\text{in particular}\qquad
 e^{X+cY}=e^Xe^{c\phi(s)Y}.
\end{equation}
Consequently, if $s\notin2\pi\ii\Z\setminus\{0\}$,
\begin{equation}\label{eq:eigenbch}
 \boxed{\displaystyle e^Xe^Y=\exp\Bigl(X+\frac{s}{1-e^{-s}}\,Y\Bigr)
 =\exp\bigl(X+\beta(s)Y\bigr),}
\end{equation}
with the coefficient interpreted as $\beta(0)=1$ at $s=0$. For every $s$,
including the resonant values,
\begin{equation}\label{eq:eigenbraid}
 e^Xe^Ye^{-X}=e^{e^sY},\qquad e^Xe^Y=e^{e^sY}e^X.
\end{equation}
The formal BCH series is
\begin{equation}\label{eq:eigenseries}
 Z(X,Y)=X+\sum_{n\geq0}\bplus{n}s^nY=X+\beta(s)Y,
\end{equation}
where the scalar series converges for $|s|<2\pi$ and diverges for
$|s|\geq2\pi$ if $Y\neq0$. All exponential identities are global in the
same settings as \cref{thm:central}.
\end{theorem}
\begin{proof}
The relation gives $\ad_X^nY=s^nY$ by induction, so the Campbell series
\eqref{eq:campbell} yields $e^{tX}Ye^{-tX}=e^{st}Y$; conjugating the
exponential series of $Y$ gives $e^{tX}e^{Y}e^{-tX}=e^{e^{st}Y}$, and at
$t=1$ this is \eqref{eq:eigenbraid}, whose second form follows by
multiplying on the right by $e^X$. No division has occurred, so these hold
for every $s$.

For \eqref{eq:eigenfactor}, let $F(t)=e^{tX}e^{c\,q_s(t)Y}$. Its right
logarithmic derivative is
\[
 F'F^{-1}=X+e^{tX}\bigl(c\,q_s'(t)Y\bigr)e^{-tX}=X+c\,e^{-st}e^{st}Y=X+cY,
\]
using $q_s'(t)=e^{-st}$ and the fact that $c\,q_s(t)Y$ commutes with its
own derivative, so that \eqref{eq:rightdexp} for the second factor reduces
to plain differentiation. Hence $F'=(X+cY)F$ with $F(0)=I$. The function
$e^{t(X+cY)}$ solves the same initial value problem, and uniqueness (as in
the proof of \cref{thm:central}) gives \eqref{eq:eigenfactor}. At $t=1$,
$q_s(1)=\phi(s)$, with the removable value $1$ at $s=0$.

If $s\notin2\pi\ii\Z\setminus\{0\}$, then $\phi(s)\neq0$ (the zeros of
$\phi$ are exactly the nonzero points of $2\pi\ii\Z$), and choosing
$c=1/\phi(s)=\beta(s)$ in \eqref{eq:eigenfactor} gives \eqref{eq:eigenbch}.

For the formal assertion \eqref{eq:eigenseries}, we show that every Lie
monomial in $X,Y$ containing at least two occurrences of $Y$ vanishes under
the relation. Consider a bracket tree evaluating such a monomial. A branch
containing no $Y$ is a bracket of copies of $X$ alone, which is $X$ if it
has length one and $0$ otherwise. A branch containing exactly one $Y$ is,
by induction on its size, a scalar multiple of $Y$: at each node one
combines either $X$ with a multiple of $Y$ (giving a multiple of $sY$) or a
zero branch with something (giving $0$). At the first node where two
branches each containing at least one $Y$ meet, both are multiples of $Y$
and their bracket is zero. Hence only the terms linear in $Y$ survive, and
by \cref{cor:linearY} these are $\beta(\ad_X)Y=\sum_n\bplus n\ad_X^nY
=\sum_n\bplus ns^nY$. The scalar series is the Taylor series of $\beta$ at
$s$, with radius exactly $2\pi$ by \cref{prop:bernoulli}; by
\eqref{eq:bernzeta} its even terms have magnitude
$2\zeta(2m)(|s|/2\pi)^{2m}\geq2(|s|/2\pi)^{2m}$, which does not tend to
zero when $|s|\geq2\pi$, so the series diverges there.
\end{proof}
The differential-equation proof of \eqref{eq:eigenbch} establishes the
exponential identity even where the Taylor series \eqref{eq:eigenseries}
diverges: resummation and Taylor convergence are distinct assertions. In
the same vein,
\begin{equation}\label{eq:eigenscaled}
 \BCH(tX,tY)=tX+t\beta(ts)Y
\end{equation}
as a Taylor series in $t$ at $0$, whose radius is exactly $2\pi/|s|$ when
$s\neq0$ and $Y\neq0$; real $s>2\pi$ therefore makes this series divergent
at $t=1$ although the resummed identity remains valid. The braiding law
\eqref{eq:eigenbraid}, being an entire exponential identity, has no
resonance exclusion at all.

\begin{example}[Resonance]\label{ex:resonance}
Let
\begin{equation}\label{eq:eigenmatrices}
 X=\begin{pmatrix}s&0\\0&0\end{pmatrix},\qquad
 Y=\begin{pmatrix}0&1\\0&0\end{pmatrix},\qquad [X,Y]=sY.
\end{equation}
Direct exponentiation gives
\[
 e^X=\begin{pmatrix}e^s&0\\0&1\end{pmatrix},\quad
 e^Xe^Y=\begin{pmatrix}e^s&e^s\\0&1\end{pmatrix},\quad
 e^{X+cY}=\begin{pmatrix}e^s&c\,(e^s-1)/s\\0&1\end{pmatrix},
\]
the last by the triangular structure ($X+cY$ has eigenvalues $s,0$ and its
exponential's off-diagonal entry is $c$ times the divided difference of
$e^z$ at $s,0$), with the continuous value $c$ at $s=0$. If $s=3\pi$, the
choice $c=\beta(3\pi)$ works although the BCH series \eqref{eq:eigenseries}
diverges. If $s=2\pi\ii k\neq0$, then $e^s=1$, every $e^{X+cY}$ equals
$I$, while $e^Xe^Y=I+Y\neq I$; thus no logarithm of the form $X+cY$ exists
at a resonance, although the logarithm $Y$ does exist, and the braiding
identity $e^Xe^Y=e^{e^sY}e^X=e^Ye^X$ remains true. What fails at resonance
is the shape $X+cY$ of the logarithm, not the existence of a logarithm.
\end{example}
